\chapter{System Design}
\label{ch:design}

\section{Requirements and Specifications}

The system must read plates in fog, rain and low contrast, work in real time as
vehicles pass, handle Indian plate formats, and run on an embedded node at each
camera. It is evaluated on recognition accuracy (clear and degraded
conditions), the effectiveness of the restoration stage, the accuracy of
multi-node correlation and the end-to-end alert latency. Table~\ref{tab:specs}
lists the design values.

\begin{table}[t]
\centering
\caption{System specifications (design values).}
\label{tab:specs}
\small
\begin{tabularx}{\textwidth}{@{}lX@{}}
\toprule
\textbf{Parameter} & \textbf{Design value}\\
\midrule
Pipeline & Capture $\rightarrow$ plate detection $\rightarrow$ weather restoration $\rightarrow$ plate reading $\rightarrow$ verification and alerts\\
Plate detector & YOLOv11n, trained on 15\,000+ annotated licence-plate images\\
Weather restoration & CLAHE for low contrast; UVMNet with physics-guided fog removal for fog; lightweight U-Net deraining for rain\\
Plate reader & LPRNet with CTC, $200\times24$ input, characters A--Z and 0--9, Indian plate-format rules\\
Confirmation & Tracking plus voting over frames: announced after 5 or more agreeing reads, or on exit for short passes\\
Edge hardware & Jetson Nano-class module, HDR / STARVIS-class camera, NoIR plate camera with IR illuminator\\
Records and alerts & Central vehicle database and multi-node correlation (Phase 2)\\
Development platform & Python and PyTorch on WSL2, NVIDIA RTX 4060 laptop GPU (8\,GB)\\
\bottomrule
\end{tabularx}
\end{table}

\section{System Architecture}

The system has two layers (Figure~\ref{fig:arch}). The \emph{camera node}
captures video and turns it into confirmed registration numbers locally; only
the recognised results, with a crop of the plate and a timestamp, are sent
onwards. The \emph{central system} verifies each number against the vehicle
database, raises alerts, forwards urgent cases to dispatch, and correlates
sightings from all nodes. Phase~1 covers the camera node.

\begin{figure}[t]
\centering
\begin{tikzpicture}[node distance=7mm and 7mm]
  % camera node
  \node[blk] (cam) {Camera\\capture};
  \node[blk,right=of cam] (det) {Plate detection\\YOLOv11n};
  \node[blk,right=of det] (res) {Weather\\restoration};
  \node[blk,right=of res] (rd) {Plate reading\\LPRNet + rules};
  \node[blk,below=9mm of rd] (trk) {Tracking and\\voting};
  \draw[arr] (cam) -- (det);
  \draw[arr] (det) -- (res);
  \draw[arr] (res) -- (rd);
  \draw[arr] (rd) -- (trk);
  \begin{scope}[on background layer]
    \node[draw=blk,dashed,rounded corners=3pt,fit=(cam)(det)(res)(rd)(trk),
          inner sep=7pt,label={[lbl]above left:CAMERA NODE (Phase 1)}] (nodebox) {};
  \end{scope}
  % central system
  \node[org,below=34mm of cam] (db) {Vehicle\\database};
  \node[org,right=of db] (cls) {Status\\classification};
  \node[org,right=of cls] (dash) {Dashboard\\and alerts};
  \node[org,right=of dash] (disp) {Dispatch\\module};
  \node[org,below=7mm of cls,minimum width=36mm] (corr) {Multi-node\\correlation};
  \draw[arr] (trk.south) -- ++(0,-6mm) -| ([xshift=4mm]cls.north)
        node[lbl,pos=0.2,above,xshift=9mm] {plate, crop, time};
  \draw[arr] (db) -- (cls);
  \draw[arr] (cls) -- (dash);
  \draw[arr] (dash) -- (disp);
  \draw[arr] (cls) -- (corr);
  \begin{scope}[on background layer]
    \node[draw=orange!70!black,dashed,rounded corners=3pt,
          fit=(db)(cls)(dash)(disp)(corr),inner sep=7pt,
          label={[lbl]below left:CENTRAL SYSTEM (Phase 2)}] {};
  \end{scope}
\end{tikzpicture}
\caption{Architecture of the Weather Adaptive ANPR System. The camera node is
implemented and tested in Phase~1; the central system is planned for Phase~2.}
\label{fig:arch}
\end{figure}

\section{Camera Node and Hardware}

Each node has an HDR or STARVIS-class camera for the scene and a NoIR camera
with an infrared illuminator for plates at night. Detection, restoration and
recognition run locally on a Jetson Nano-class module with SSD storage in a
weatherproof enclosure. Local processing keeps the network load small, because
only recognised results are sent to the central system, and keeps the node
working if the link is lost.

The platform was chosen with the restoration stage in mind. The first design
targeted a Raspberry Pi~5, but the Mamba selective-scan kernel used by UVMNet
is CUDA-only and does not run there; a transformer-only restorer was
therefore built for the Pi and later replaced by UVMNet once a CUDA-capable
embedded target was chosen. Whether the final network meets the latency budget
on the Jetson module is one of the Phase~2 tests.

\section{Software Pipeline}

The pipeline is written in Python with PyTorch and the Ultralytics YOLO
package, OpenCV for image handling, and runs on WSL2 with an NVIDIA RTX~4060
laptop GPU during development. Figure~\ref{fig:pipeline} shows the processing
of one frame.

\begin{figure}[t]
\centering
\begin{tikzpicture}[node distance=6mm and 6mm]
  \node[blk] (frame) {Frame};
  \node[blk,right=of frame] (route) {Weather\\check};
  \node[grn,above right=4mm and 9mm of route] (clahe) {CLAHE\\(low contrast)};
  \node[grn,right=9mm of route] (uvm) {UVMNet +\\physics (fog)};
  \node[grn,below right=4mm and 9mm of route] (rain) {Deraining\\(rain)};
  \node[blk,right=44mm of route] (yolo) {YOLOv11n\\plate boxes};
  \node[blk,below=8mm of yolo] (crop) {Crop,\\flatten 2-line};
  \node[blk,left=of crop] (lpr) {LPRNet\\+ format decode};
  \node[blk,left=of lpr] (trk) {Track and\\vote};
  \node[blk,left=of trk] (out) {Confirmed\\plate};
  \draw[arr] (frame) -- (route);
  \draw[arr] (route.east) -- ++(4mm,0) |- (clahe.west);
  \draw[arr] (route.east) -- (uvm.west);
  \draw[arr] (route.east) -- ++(4mm,0) |- (rain.west);
  \draw[arr] (clahe.east) -- ++(5mm,0) |- (yolo.west);
  \draw[arr] (uvm.east) -- (yolo.west);
  \draw[arr] (rain.east) -- ++(5mm,0) |- (yolo.west);
  \draw[arr] (yolo) -- (crop);
  \draw[arr] (crop) -- (lpr);
  \draw[arr] (lpr) -- (trk);
  \draw[arr] (trk) -- (out);
\end{tikzpicture}
\caption{Processing of one frame in the camera node. The weather check selects
the restoration route; plate detection, reading and voting follow. The order of
restoration and detection can be swapped so that only the plate crop is
restored.}
\label{fig:pipeline}
\end{figure}

Two orders of restoration and detection are supported. In the first the whole
frame is restored before detection, which gives the clearest picture and is
used for the foggy-video pipeline of Section~\ref{sec:foggyvideo}. In the
second the plate is detected first and only its crop is restored and read,
which is cheaper and is the order of the real-time loop. The reader is given a
``select'' view that chooses between the plain crop and a smoothed crop, because
experiments (Section~\ref{sec:ponmudi}) showed that heavy restoration of a
crop can hurt reading.

\section{Data Handling}

Detections are tracked from frame to frame, and the readings of each track are
kept in a vote table. A reading is accepted only if it matches an Indian plate
layout, and a plate is confirmed when at least five reads agree or, for a
vehicle that is in view only briefly, announced when it leaves the scene if the
reading is well formed and most reads agree
(Section~\ref{sec:realtime}). The central system receives the confirmed plate
number, a crop of the plate, the camera identifier and the time.

\section{Summary}

The design separates a camera node, which turns pixels into confirmed plate
numbers on embedded hardware, from a central system that verifies, alerts and
correlates. The weather-dependent part is isolated in a restoration stage that
is selected by a weather check, so that clear frames are not processed
unnecessarily and each degradation has its own model.
